% SGA 3, Expose IV, Corollary 4.5.4 through Scholium 4.5.5.
% French authority: Exp4-13oct24.pdf, local pages 45--46
% (combined-reader pages 233--234; printed pages 221--222).
% The external English Markdown was used only as a comparison.

\medskip
\noindent\SGARefTarget{sga3:iv:corollary:4:5:4}\textbf{推论 4.5.4.}
\textit{假设 \(\mathscr C\) 上给定的拓扑比典范拓扑更粗.
对每个 \(S\in\operatorname{Ob}(\mathscr C)\), 都有范畴等价}
\[
  \widetilde{\mathscr C}_{/\widetilde S}
  \xrightarrow{\ \sim\ }
  \widetilde{\mathscr C/S}.
\]
\textit{下列各图在同构意义下交换;
图中每个未标记的箭头都是等价:}
\[
\begin{array}{c@{\;}c@{\;}c@{\;}c@{\;}c@{\;}c@{\;}c}
\mathscr C_{/S}
 &\xrightarrow{(j_{\mathscr C})_{/S}}&
 \widetilde{\mathscr C}_{/\widetilde S}
 &\xrightarrow{(i_{\mathscr C})_{/\widetilde S}}&
 \widehat{\mathscr C}_{/\widehat S}
 &\xrightarrow{(a_{\mathscr C})_{/\widehat S}}&
 \widetilde{\mathscr C}_{/\widetilde S}
 \\
 \big\Vert && \big\downarrow\wr && \big\downarrow\wr && \big\downarrow\wr
 \\
\mathscr C_{/S}
 &\xrightarrow{j_{\mathscr C/S}}&
 \widetilde{\mathscr C/S}
 &\xrightarrow{i_{\mathscr C/S}}&
 \widehat{\mathscr C/S}
 &\xrightarrow{a_{\mathscr C/S}}&
 \widetilde{\mathscr C/S}.
\end{array}
\]
若 \(T\to S\) 是 \(\mathscr C/S\) 的一个对象, 则还有交换图
\[
\begin{array}{ccccc}
 \bigl(\widetilde{\mathscr C}_{/\widetilde S}\bigr)_{/\widetilde T}
 &\longrightarrow&
 \bigl(\widetilde{\mathscr C/S}\bigr)_{/\widetilde T}
 &&\\[-0.15em]
 \Big\downarrow\rlap{\(\scriptstyle\wr\)}
 &&& \searrow &
 \widetilde{(\mathscr C/S)_{/T}}\, .
 \\
 \widetilde{\mathscr C}_{/\widetilde T}
 &\longrightarrow&
 \widetilde{\mathscr C/T}
 &\mathrel{\overset{\sim}{\nearrow}}&
\end{array}
\]

前两个方块的交换性, 由等价
\(
  \widetilde{\mathscr C}_{/\widetilde S}
  \to\widetilde{\mathscr C/S}
\)
的定义得到.
要证明最后一个方块的交换性, 只需注意下列方块交换:
\[
\begin{array}{ccc}
 \widehat{\mathscr C}_{/\widehat S}
 &\xrightarrow{L'}&
 \widehat{\mathscr C}_{/\widehat S}
 \\
 \big\downarrow && \big\downarrow
 \\
 \widehat{\mathscr C/S}
 &\xrightarrow{L''}&
 \widehat{\mathscr C/S}.
\end{array}
\]
这里, \(L'\) 是函子 \(L_{\mathscr C}\) 在
\(\widehat{\mathscr C}_{/\widehat S}\) 上的限制,
而 \(L''\) 是函子 \(L_{\mathscr C/S}\).
回到 \SGARefLink{sga3:iv:corollary:4:3:9}{4.3.9} 之后给出的函子 \(L\)
的\SGARefLink{sga3:iv:definition:4:3:10:0}{定义}, 立即可以看出这一点.
第二个图只是预层范畴的相应交换图 (第 I 篇,
\SGARefLink{sga3:i:section:1}{第 1 号}) 在层范畴上的限制.

\medskip
\noindent\SGARefTarget{sga3:iv:scholium:4:5:5}\textbf{附论 4.5.5.}
以上断言表明, 当给定的拓扑比典范拓扑更粗时,
可以将 \(\mathscr C\) 等同于 \(\widetilde{\mathscr C}\) 的一个全子范畴,
后者本身又是 \(\widehat{\mathscr C}\) 的一个全子范畴.
在这些等同下, 可以对 \(\mathscr C\to
\widehat{\mathscr C}\) 作通常的语言滥用,
以上交换图保证了这些用法的合理性.
特别地, \SGARefLink{sga3:iv:corollary:4:5:4}{4.5.4} 的第一个图表明,
使用函子 \(a\) 时无需特别留意.

下一小节将说明, 将 \(\mathscr C\) 等同于
\(\widetilde{\mathscr C}\) 的全子范畴与某些正向极限的构造可交换,
这有别于将其等同于 \(\widehat{\mathscr C}\) 的全子范畴的情形;
随后将说明如何运用这一事实.

\emph{从现在起, 除非明确另作声明,
都假设给定的拓扑比典范拓扑更粗,
并系统地作上述等同.}

% Translation coverage: entire assigned source component, without omissions.
% Statistics: source lines=106; target lines=105; inline math=20; display math=4; unresolved translation TODOs=0.
